% Source: https://artofproblemsolving.com/community/c3h1210385_data_points
% Stable identity: aops-topic-1210385-q01
Only one value is not prime. Each mode occurs twice, so both modes are prime, and they differ by $2$: the modes are $3,5$ or $5,7$. The other three values are distinct and differ from the modes, and exactly one of them is not prime.
\begin{itemize}
\item Modes $3,5$: the two other primes must be $2$ and $7$. The range is then $7-2=5$, so the non-prime $n$ lies between $2$ and $7$: $n=4$ or $6$. The data $2,3,3,4,5,5,7$ has median $4$ (not a mode), while $2,3,3,5,5,6,7$ has median $5$. \checkmark
\item Modes $5,7$: the two other primes must be $2$ and $3$; again the range is $5$ and $n=4$ or $6$. Both $2,3,4,5,5,7,7$ and $2,3,5,5,6,7,7$ have median $5$. \checkmark
\end{itemize}
So there are $\boxed{3}$ data sets: $\{2,3,3,5,5,6,7\}$, $\{2,3,4,5,5,7,7\}$ and $\{2,3,5,5,6,7,7\}$.
